\documentclass[applsci,article,submit,moreauthors]{Definitions/mdpi}

\firstpage{1}
\pubvolume{1}
\issuenum{1}
\articlenumber{0}
\pubyear{2026}
\copyrightyear{2026}

% Keep the MDPI layout while identifying this version as a manuscript draft.
\let\linenumbers\nolinenumbers
\setlength{\headheight}{21pt}
\addtolength{\topmargin}{-9pt}
\newcommand{\DraftStatus}{Manuscript draft for \textit{Applied Sciences}}
\fancyhead[L]{\footnotesize\vspace{2pt}\DraftStatus\\\rule{\fulllength}{0.4pt}}
\fancyfoot[R]{}
\makeatletter
\apptocmd{\ps@plain}{\fancyfoot[L]{\footnotesize\DraftStatus}\fancyfoot[R]{}}{}{}
\makeatother
\renewcommand{\cright}{\DraftStatus.}

\newcommand{\softwarepm}{PM4Py}
\newlength{\PlaceholderSide}
\newcommand{\FigurePlaceholder}[2][0.35]{%
  \begingroup
  \setlength{\PlaceholderSide}{#1\linewidth}%
  \fbox{\parbox[c][\PlaceholderSide][c]{\PlaceholderSide}{%
    \centering\small\hspace{0pt}#2\par}}%
  \endgroup
}

\Title{Detecting Concept Drift with Process States in Case-Centric Event Logs}
\Author{Bora Deren and Alessandro Berti *}
\AuthorNames{Bora Deren and Alessandro Berti}
\address{Chair of Process and Data Science, RWTH Aachen University, Aachen, Germany}
\corres{Correspondence: a.berti@pads.rwth-aachen.de}

\abstract{A process can change who does the work or how long cases wait while keeping the same sequence of activities. We present Kairo, an approach that learns recurring process states from case-centric event logs. It describes progress within cases, patterns of resource use, and conditions across cases with separate sets of features. It groups similar observations into states and compares how common those states are over time. A graphical interface lets analysts inspect the states and the features behind a drift signal. In a retrospective benchmark, Kairo located selected changes in resource assignment and waiting time without extra alarms. Its strongest overall configuration matched more annotated changes than the strongest comparison configuration and issued more unmatched alarms. The results support using process states to investigate change and motivate further work on the reliability of automatic alerts.}

\keyword{process mining; concept drift; event logs; process monitoring; clustering; visual analytics}

\begin{document}
\renewcommand{\sectionautorefname}{Section}

\section{Introduction}
\label{sec:introduction}

Process mining uses records of work to understand how a process operates~\cite{vanderaalst2016process}. A model discovered from an entire event log can combine behavior from periods in which the process operated differently. A change in the underlying behavior is called \emph{concept drift}. Identifying when a change occurred helps an analyst choose which periods to compare. Understanding what changed helps the analyst decide which cases, activities, or resources to examine.

A change in activity order is visible to a detector that compares sequences of activities. Other changes require different observations. For example, a service desk may route the same requests through the same activities while reassigning the work to another team. An increase in incoming requests may also cause more cases to be open at the same time, even if each request follows its usual path. Research has emphasized changes in activity order, although other approaches track time, resources, or case attributes~\cite{sato2022survey}. The practical gap we address is how to summarize several related measurements within each perspective and let analysts inspect how those summaries change over time.

We present Kairo, which builds on the state-based monitoring method developed in Deren's bachelor's thesis~\cite{kairo-thesis}. Kairo represents the progress of individual cases, the work carried out by resources, and the conditions shared by cases as they run. Each perspective uses features suited to what it observes. Clustering turns groups of similar observations into recurring states. A change in the mix of states produces a drift signal, and the measurements behind the states give the analyst a starting point for interpreting it. A Python library and graphical interface support this analysis.

This paper explains the method and compares its drift estimates with those of existing detectors on logs with known changes and on controls without injected changes. The comparison uses automatic alarm settings selected without consulting the drift annotations. The results show that the state signals can surface changes in resource assignment and waiting time. They also show that greater coverage of annotated changes can come with more alarms that do not match an annotation.

The rest of the paper is organized as follows. \autoref{sec:related-work} reviews existing drift detection approaches. \autoref{sec:preliminaries} introduces the terms used in \autoref{sec:approach}, which presents the method. \autoref{sec:implementation} describes the library and interface. \autoref{sec:assessment} presents the assessment, and \autoref{sec:discussion} discusses its findings and limits. \autoref{sec:conclusion} concludes the paper and outlines future work.

\section{Related Work}
\label{sec:related-work}

Research on concept drift in process mining has considered several ways to describe and compare process behavior over time. Control-flow approaches use relations between activities or distributions of execution traces. Other approaches track timing, resource use, and case characteristics. Changes can also be detected by monitoring conformance to process models or constraints, by recognizing patterns learned from labeled examples, or by grouping similar observations and tracking shifts in those groups.

\subsection{Activity Relations and Trace Distributions}

Bose et al. characterize control flow through relations between activities. Their relation type count records how many activities always, sometimes, or never follow each activity in a log fragment \cite{bose2011handling}. ProDrift tests whether the frequencies of partially ordered runs change across neighboring windows. By grouping traces that differ only in the order of concurrent activities, it avoids treating every observed ordering as a different run \cite{maaradji2017detecting}. The event-stream method of Ostovar et al. compares counts of relations extracted from events, allowing it to produce evidence while cases are still open \cite{ostovar2016streams}.

Local Completeness-based Drift Detection looks for relations that appear or cease to occur, while relation invariance uses changes in the presence of directly-follows and eventually-follows relations \cite{lin2022lcdd,zheng2017detecting}. Graph metrics summarize the structure of a process graph, for example through branching and connectivity \cite{seeliger2017graph}. In Kairo, activity and relation features help describe the observed prefix of a running case. Similar prefixes can share a state, and drift is assessed through changes in the frequencies of those states. The resource and inter-case analyses use other measurements and apply the same comparison to their state frequencies.

\subsection{Drift beyond Control Flow}

Brockhoff et al. use the Earth Mover's Distance to compare distributions of traces. The distance accounts for how similar traces are, and the time-aware representation includes service or sojourn times \cite{brockhoff2020time}. This retains differences between traces that can be lost when similar observations are assigned to the same state. Kairo can also represent conditions such as workload and resource allocation across cases in a calendar window.

Adams et al. detect and relate changes in feature series covering multiple process perspectives \cite{adams2023explainable}. Klijn et al. extend multi-perspective drift detection to actor behavior and collaboration using an event knowledge graph \cite{klijn2024multi}. Khraiwesh et al. use the coefficient of variation to choose a smoothing width for each indicator before detecting change points \cite{khraiwesh2027window}. Multivariate tests such as maximum mean discrepancy can detect changes in the joint distribution of feature vectors \cite{gretton2012kernel}. Kairo groups related measurements into state descriptions for each perspective, such as a workload state that combines arrivals with the cases already in progress. The resulting state can recur, and its share can be examined alongside the measurements typical of that state.

\subsection{Conformance, Constraints, and Supervised Detection}

FAIRY monitors how incoming traces conform to a reference model and updates the model after a detected change \cite{pasquadibisceglie2026fairy}. CRIER tracks fitness and precision over windows. Precision can help reveal a change when cases still fit the model but some allowed behavior disappears from the observations \cite{gallego2025crier}. Visual Drift Detection examines changes in declarative constraints, such as how regularly one activity is followed by another \cite{yeshchenko2019visual}. These methods provide evidence in terms of the model or rules they monitor. Kairo learns state descriptions directly from the selected log features and uses a reference distribution of states.

Kraus and van der Aa use labeled examples to train a detector to recognize drift patterns in images of window similarities \cite{kraus2025machine}. Kairo learns its states without drift annotations. Its divergence signal indicates when state usage changes and can guide inspection of the underlying observations.

\subsection{Clustering and Process States}

Clustering provides a way to summarize behavior before looking for change. Barbon Junior et al. cluster traces in an event stream and use the emergence of a new cluster as a signal \cite{barbonjunior2018framework}. Version Clustering groups windows into process versions based on similarities in their directly-follows profiles and estimates boundaries between versions \cite{abarcazuniga2026version}. Kairo uses clustering to summarize event prefixes or calendar windows, depending on the perspective. Several states can be present throughout a stable period, and a change may be visible as a shift in their relative frequency.

The state representation builds on the work of Berti et al., who use a self-organizing map to learn execution states from sequences of events associated with a business object \cite{berti2025ocel}. Kairo uses case prefixes to describe within-case behavior and extends the state representation to resource and inter-case conditions over calendar time. Steady-state detection by Kraus et al. identifies periods of stable process behavior for tasks such as discovery and performance analysis \cite{kraus2026steady}. Kairo uses states as recurring descriptions of observed behavior and tracks how their distribution changes over time.

\section{Preliminaries}\label{sec:preliminaries}

A case is an instance of a process, such as a request handled by a service team. In a traditional event log, each event records an activity, a timestamp, and the case to which it belongs. An event can also record who performed the activity and when it started. Case attributes, such as the category of a request, describe the case as a whole. A \emph{trace} is the sequence of observed events for a case in timestamp order, with ties resolved using their order in the input. The \emph{prefix} at an event contains that event and its predecessors in the same trace. A \emph{directly-follows relation} is an ordered pair of activities at consecutive positions in a trace.

A \emph{calendar window} is a fixed-length interval on a time grid. It includes its start and excludes its end. The grid lets us group events from different cases by time, even when a case spans several windows. An event at a boundary belongs to the window that begins there.

In a sudden drift, the change takes place at a boundary. In a gradual drift, the change unfolds over a period. A previous pattern can also recur. We use \emph{state} to mean a cluster of similar observations. In this sense, a state describes a local pattern of behavior that can occur repeatedly within the same process regime. Drift evidence comes from changes in how often these states occur over time.

\section{Approach}\label{sec:approach}

Kairo constructs a feature matrix for each perspective, scales and optionally reduces it, assigns states by clustering, and compares state distributions over time. \autoref{fig:approach-overview} shows how the stages connect.

\subsection{Representing Process Behavior}\label{sec:approach-features}

The unit of observation depends on the perspective. Event-level observations retain position within a case. Window-level observations summarize conditions shared across cases over a calendar interval. \Cref{def:feature-representation} defines these observation sequences and the feature rows attached to them.

\begin{Definition}[Observation sequences and feature maps]\label{def:feature-representation}
Write $p$ for the selected perspective, with $p\in\{\mathrm{intra},\mathrm{res},\mathrm{inter}\}$.
\begin{itemize}[leftmargin=*,itemsep=.2em,topsep=.3em]
  \item \textbf{Intra-case observations.} The observation sequence is $\sigma_{\mathrm{intra}}=\sigma(E)$, with one observation per event and time anchor $t_{\mathrm{intra}}(e)=t(e)$. Its feature history comes from the trace of its own case.
  \item \textbf{Resource and inter-case observations.} Choose a width $\delta_p$ and use the covering grid segment $\sigma_p=\mathcal W_{\delta_p}[\sigma(E)]$ from \Cref{def:calendar-grid}. Each interval $W$ in this sequence gives one feature window, with anchor $t_p(W)=b(W)$. These sequences retain intervals without events.
  \item \textbf{Feature map.} Let $O_p=\operatorname{set}(\sigma_p)$ be the set of observations. The map $x_p\colon O_p\to\mathbb R^{d_p}$ assigns each observation a vector containing its $d_p$ selected feature coordinates.
  \item \textbf{Feature values.} For a proposition $B$, the indicator $\mathbf{1}[B]$ is one if $B$ holds and zero otherwise. Ratios with zero denominator, means of empty samples, and sample standard deviations based on fewer than two values are set to zero. Elapsed gaps and event durations are expressed in minutes.
\end{itemize}
\end{Definition}

The time anchor links each row to the sequence used for drift scoring in \Cref{sec:approach-signal}. The following definitions specify which events contribute to the feature coordinates.

\subsubsection{Intra-Case Perspective}

\Cref{def:intra-features} describes the current activity, earlier activities, and how far the case has progressed. Each event is evaluated in the order of its own trace.

\begin{Definition}[Intra-case features]\label{def:intra-features}
For $c_0\in C$, use the event subset $E_{c_0}$ and its ordered trace $\tau_{c_0}$ from \Cref{def:ordered-log}. For $e=e_{c_0,j}$, the feature history is the position window $\pi(e)=\tau_{c_0}[1:j]$ from \Cref{def:case-history}. For $a^\star\in A$, let $N_{a^\star}(e)=|\{v\in\operatorname{set}(\pi(e)):a(v)=a^\star\}|$.
\begin{itemize}[leftmargin=*,itemsep=.2em,topsep=.3em]
  \item \textbf{Current activity.} The coordinate for $a^\star$ is $\mathbf{1}[a(e)=a^\star]$.
  \item \textbf{Past activities.} For a selected integer lag limit $h\geq0$, use one indicator block for each $\ell\in\{1,\ldots,h\}$. The coordinate for $a^\star$ in block $\ell$ is $\mathbf{1}[a(e_{c_0,j-\ell})=a^\star]$ if $j>\ell$, and zero otherwise.
  \item \textbf{Prefix presence.} The indicator for $a^\star$ is $\mathbf{1}[N_{a^\star}(e)>0]$.
  \item \textbf{Prefix frequency.} The share of $a^\star$ is $N_{a^\star}(e)/j$.
  \item \textbf{Directly-follows counts.} For each pair $(a^\star,a')$ observed in the log, the number of its occurrences within the prefix is
  \[
    \bigl|\{v\in\operatorname{set}(\pi(e)):v^-\text{ exists},\
                 (a(v^-),a(v))=(a^\star,a')\}\bigr|.
  \]
  \item \textbf{Case progress.} The normalized position is $j/m_{c_0}$.
\end{itemize}
The selected groups of these coordinates form $x_{\mathrm{intra}}(e)$.
\end{Definition}

The prefix ends at the current event. The optional progress coordinate uses the complete recorded trace length, so it provides a retrospective measure of position.

\subsubsection{Resource Perspective}

\Cref{def:resource-features} groups resource activity into the feature windows from \Cref{def:feature-representation}. It assigns a handover to the interval containing the receiving event, which allows transfers across window boundaries to contribute.

\begin{Definition}[Resource features]\label{def:resource-features}
Let $E^{\mathrm{res}}=\{e\in E:r(e)\ne\bot\}$ be the events with recorded resources. For $W\in O_{\mathrm{res}}$, apply \Cref{def:calendar-grid} to $\sigma(E^{\mathrm{res}})$ and set $E_W^{\mathrm{res}}=\operatorname{set}(\sigma(E^{\mathrm{res}})_W)$. For $\rho\in R$, let $E_{W,\rho}=\{e\in E_W^{\mathrm{res}}:r(e)=\rho\}$. The receiving events of handovers are
\[
 H_W=\{e\in E_W^{\mathrm{res}}: e^-\text{ exists},\
                                  r(e^-)\notin\{\bot,r(e)\}\}.
\]
The predecessor is taken from the full case trace. For $\rho,\nu\in R$, let $H_{W,\rho,\nu}=\{e\in H_W:r(e^-)=\rho,\ r(e)=\nu\}$.
\begin{itemize}[leftmargin=*,itemsep=.2em,topsep=.3em]
  \item \textbf{Events per resource.} The count for $\rho$ is $|E_{W,\rho}|$.
  \item \textbf{Cases per resource.} The distinct case count is $|\{c(e):e\in E_{W,\rho}\}|$.
  \item \textbf{Mean event duration.} Average the available $\operatorname{dur}(e)$ over $e\in E_{W,\rho}$.
  \item \textbf{Mean incoming handover gap.} Average $g(e)$ over $e\in H_W\cap E_{W,\rho}$.
  \item \textbf{Activity-resource share.} For each observed pair $(a^\star,\rho)$, the share of that activity performed by $\rho$ is
  \[
     \frac{|\{e\in E_{W,\rho}:a(e)=a^\star\}|}
          {|\{e\in E_W^{\mathrm{res}}:a(e)=a^\star\}|}.
  \]
  \item \textbf{Handover share.} For each observed handover pair $(\rho,\nu)$, the share of transfers from $\rho$ received by $\nu$ is
  \[
     \frac{|H_{W,\rho,\nu}|}{\sum_{\kappa\in R}|H_{W,\rho,\kappa}|}.
  \]
\end{itemize}
The selected feature groups and resources determine $x_{\mathrm{res}}(W)$. Each retained handover-share coordinate uses two selected resources. The denominators in both shares use all recorded resources.
\end{Definition}

Activity shares can reveal a reassignment even when the total number of events stays stable. The full denominators preserve that interpretation when the analyst displays a subset of resources.

\subsubsection{Inter-Case Perspective}

\Cref{def:inter-features} uses the full event sequence to summarize activity and timing across cases. The distinction between case arrivals and completions helps separate changes in incoming demand from changes in how quickly work finishes.

\begin{Definition}[Inter-case features]\label{def:inter-features}
For $W\in O_{\mathrm{inter}}$, take the event set $E_W$ obtained by applying \Cref{def:calendar-grid} to the full sequence $\sigma(E)$. The cases observed in the window form $C_W=\{c(e):e\in E_W\}$.
\begin{itemize}[leftmargin=*,itemsep=.2em,topsep=.3em]
  \item \textbf{Active cases.} The count is $|C_W|$.
  \item \textbf{Arrivals and completions.} Count cases with first or last recorded timestamps in the interval, respectively. The counts are $|\{c_0\in C:s_{c_0}\in W\}|$ and $|\{c_0\in C:f_{c_0}\in W\}|$.
  \item \textbf{Events per active case.} The ratio is $|E_W|/|C_W|$.
  \item \textbf{Within-case gaps.} Take the mean and sample standard deviation of $(g(e):e\in E_W,\ e^-\text{ exists})$. A gap contributes to the window of its later event.
  \item \textbf{Stalled cases.} Given a stall threshold $\eta_{\mathrm{stall}}>0$, the number of cases stalled at the interval end is
  \[
     \bigl|\{c_0\in C:\exists\,e\in E_{c_0}\ \text{with}\ e^+\text{ defined and }
                 t(e)+\eta_{\mathrm{stall}}<d(W)<t(e^+)\}\bigr|.
  \]
  \item \textbf{Numerical attributes.} For each mapped numerical case attribute $\alpha$, take the mean and sample standard deviation of the available values $(v_\alpha(e):e\in E_W,\ v_\alpha(e)\ne\bot)$.
  \item \textbf{Categorical attributes.} For each observed value $z$ of a mapped categorical case attribute $\alpha$, its share is $|\{e\in E_W:v_\alpha(e)=z\}|/|E_W|$.
\end{itemize}
The selected groups form $x_{\mathrm{inter}}(W)$. Attribute summaries count each case once per event it contributes.
\end{Definition}

A stalled case may have no events in the current feature window. For categorical shares, the denominator includes events whose category is missing.

\subsubsection{Running Example}\label{sec:approach-example}

Consider a service desk observed over consecutive feature windows $W_{\mathrm a},\ldots,W_{\mathrm f}$. In each window $W_{\mathrm x}$, a focal case $\mathrm x$ and its companion case overlap in time. Each has a Receive event followed by a Resolve event, and both events fall within that window. The Frontline team receives every case and resolves every companion case after a gap of $g$. It also resolves the focal cases $\mathrm a$, $\mathrm b$, and $\mathrm c$ after the same gap. The Specialist team resolves the focal cases $\mathrm d$, $\mathrm e$, and $\mathrm f$ after a longer gap $G>g$.

\begin{Example}[Intra-case features for a case prefix]\label{ex:intra-case}
Compare the Receive and Resolve prefixes of cases $\mathrm c$ and $\mathrm d$. At Receive, the current-activity indicator for Receive is on and the prefix contains no directly-follows pair. At Resolve, Receive is the immediately preceding activity and the prefix contains the Receive-to-Resolve pair once. \autoref{tab:example-intra-case} shows the same feature values at the corresponding positions of both cases. \autoref{fig:example-intra-case} illustrates the prefixes used to build the rows.
\end{Example}

\begin{table*}[!t]
  \caption{Selected intra-case features for the Receive and Resolve events of cases $\mathrm c$ and $\mathrm d$. On and off denote true and false for an indicator.}
  \label{tab:example-intra-case}
  \centering
  \small
  \begin{tabularx}{\linewidth}{@{}ll*{3}{>{\centering\arraybackslash}X}@{}}
    \toprule
    \textbf{Case} & \textbf{Current event} & \textbf{Receive now} &
      \textbf{Receive immediately before} & \textbf{Receive-to-Resolve count} \\
    \midrule
    $\mathrm c$ & Receive & On & Off & Zero \\
    $\mathrm c$ & Resolve & Off & On & One \\
    $\mathrm d$ & Receive & On & Off & Zero \\
    $\mathrm d$ & Resolve & Off & On & One \\
    \bottomrule
  \end{tabularx}
\end{table*}

\begin{figure*}[!t]
  \centering
  \begin{tikzpicture}[
  kairo/base, x=1cm, y=1cm,
  activity/.style={
    rounded corners=2pt, minimum width=1.58cm, minimum height=0.37cm,
    draw=KairoLine, fill=white, inner sep=2pt
  },
  current/.style={activity, draw=KairoBlue, fill=KairoBlueTint, text=KairoBlue},
  follows/.style={-{Stealth[length=4pt]}, draw=KairoMuted, line width=0.45pt}
]
  \path[use as bounding box] (0,0) rectangle (13.08,-1.56);
  \path[fill=KairoPaper, rounded corners=3pt] (1.69,-1.55) rectangle (7.23,-0.40);
  \path[fill=KairoPaper, rounded corners=3pt] (7.49,-1.55) rectangle (13.03,-0.40);
  \node[anchor=west, font=\footnotesize\bfseries] at (0,-0.16) {Case};
  \node[font=\footnotesize\bfseries] at (4.46,-0.16) {Prefix at Receive};
  \node[font=\footnotesize\bfseries] at (10.26,-0.16) {Prefix at Resolve};
  \foreach \caseid/\rowy in {c/-0.72,d/-1.23} {
    \node[anchor=west] at (0,\rowy) {$\mathrm{\caseid}$};
    \node[current] at (4.46,\rowy) {Receive};
    \node[activity, text=KairoMuted] (receive-\caseid) at (9.29,\rowy) {Receive};
    \node[current] (resolve-\caseid) at (11.23,\rowy) {Resolve};
    \draw[follows] (receive-\caseid.east) -- (resolve-\caseid.west);
  }
\end{tikzpicture}

  \caption{Cases $\mathrm c$ and $\mathrm d$ have matching activity prefixes at each event. The shaded activity is the current event in each prefix.}
  \label{fig:example-intra-case}
\end{figure*}

\begin{Example}[Resource features for a handover]\label{ex:resource}
Between $W_{\mathrm c}$ and $W_{\mathrm d}$, responsibility for resolving the focal case moves from Frontline to Specialist. Frontline still resolves the companion case, so Specialist handles half of the Resolve events in $W_{\mathrm d}$. The only handover in that window goes from Frontline to Specialist, with a gap of $G$. \autoref{tab:example-resource} shows how the assignment and handover features change. \autoref{fig:example-resource} shows who performs each focal event.
\end{Example}

\begin{table*}[!t]
  \caption{Selected resource features in consecutive feature windows. The handover share and mean incoming gap are zero in $W_{\mathrm c}$, which contains no handover.}
  \label{tab:example-resource}
  \centering
  \small
  \begin{tabularx}{\linewidth}{@{}X*{2}{>{\centering\arraybackslash}p{0.17\linewidth}}@{}}
    \toprule
    \textbf{Feature} & $\boldsymbol{W_{\mathrm c}}$ & $\boldsymbol{W_{\mathrm d}}$ \\
    \midrule
    Events by Frontline & Four & Three \\
    Events by Specialist & Zero & One \\
    Share of Resolve performed by Specialist & Zero & Half \\
    Share of Frontline's handovers received by Specialist & Zero & Full \\
    Mean incoming handover gap for Specialist & Zero & $G$ \\
    \bottomrule
  \end{tabularx}
\end{table*}

\begin{figure*}[!t]
  \centering
  \begin{tikzpicture}[
  kairo/base, x=1cm, y=1cm,
  receive/.style={circle, minimum size=4.3pt, inner sep=0pt, line width=0.75pt, fill=white},
  resolve/.style={rectangle, minimum size=4.1pt, inner sep=0pt},
  same team/.style={draw=KairoBlue, line width=0.85pt, line cap=round},
  handover/.style={draw=KairoAmber, line width=0.85pt, dash pattern=on 3pt off 1.5pt,
    -{Stealth[length=4pt]}, shorten >=3pt}
]
  \path[use as bounding box] (0,0) rectangle (13.08,-2.75);
  % Use the same local time positions as the inter-case diagram.
  \foreach \windowid/\leftx/\windowtitle/\focalend/\focaly in {
    c/2.10/Earlier/2.84/-1.31,
    d/7.81/Later/4.49/-1.97
  } {
    \begin{scope}[shift={(\leftx,0)}]
      \node[font=\footnotesize\bfseries] at (2.635,-0.17)
        {\windowtitle\ window $W_{\mathrm{\windowid}}$};
      \path[fill=KairoBlueTint, rounded corners=2pt] (0,-1.58) rectangle (5.27,-0.50);
      \path[fill=KairoAmberTint, rounded corners=2pt] (0,-2.26) rectangle (5.27,-1.64);
      \coordinate (companion-start-\windowid) at (1.46,-0.81);
      \coordinate (companion-end-\windowid) at (3.56,-0.81);
      \coordinate (focal-start-\windowid) at (0.74,-1.31);
      \coordinate (focal-end-\windowid) at (\focalend,\focaly);
    \end{scope}
  }
  \node[anchor=east, font=\footnotesize\bfseries] at (1.96,-0.49) {Frontline};
  \node[anchor=east, text=KairoMuted] at (1.96,-0.81) {Companion};
  \node[anchor=east, text=KairoMuted] at (1.96,-1.31) {Focal case};
  \node[anchor=east, font=\footnotesize\bfseries] at (1.96,-1.97) {Specialist};
  \foreach \windowid in {c,d} {
    \draw[same team] (companion-start-\windowid) -- (companion-end-\windowid);
    \node[receive, draw=KairoBlue] at (companion-start-\windowid) {};
    \node[resolve, fill=KairoBlue] at (companion-end-\windowid) {};
  }
  \draw[same team] (focal-start-c) -- (focal-end-c);
  \draw[handover] (focal-start-d) -- (focal-end-d);
  \node[anchor=south, text=KairoAmber, inner sep=1pt] at (10.69,-1.59) {Handover};
  \node[receive, draw=KairoBlue] at (focal-start-c) {};
  \node[receive, draw=KairoBlue] at (focal-start-d) {};
  \node[resolve, fill=KairoBlue] at (focal-end-c) {};
  \node[resolve, fill=KairoAmber] at (focal-end-d) {};
  \node[receive, draw=KairoMuted] at (4.85,-2.60) {};
  \node[anchor=west, text=KairoMuted] at (4.98,-2.60) {Receive};
  \node[resolve, fill=KairoMuted] at (6.56,-2.60) {};
  \node[anchor=west, text=KairoMuted] at (6.69,-2.60) {Resolve};
  \draw[-{Stealth[length=4pt]}, draw=KairoMuted, line width=0.5pt]
    (8.19,-2.60) -- (8.74,-2.60);
  \node[anchor=west, text=KairoMuted] at (8.84,-2.60) {Time within window};
\end{tikzpicture}

  \caption{The focal case passes from Frontline to Specialist in $W_{\mathrm d}$. The handover is assigned to the window containing the Specialist event.}
  \label{fig:example-resource}
\end{figure*}

\begin{Example}[Inter-case features for a longer gap]\label{ex:inter-case}
Both $W_{\mathrm c}$ and $W_{\mathrm d}$ have the same numbers of arrivals, completions, and active cases. Each Resolve event contributes its gap from Receive to the mean for its window. The gaps are both $g$ in $W_{\mathrm c}$. In $W_{\mathrm d}$, the focal case contributes $G$ and the companion still contributes $g$, so the mean rises. \autoref{tab:example-inter-case} shows this change in timing alongside the matching counts. \autoref{fig:example-inter-case} shows which gaps enter each window's mean.
\end{Example}

\begin{table*}[!t]
  \caption{Selected inter-case features in consecutive feature windows. The later window has the same counts and a longer mean gap.}
  \label{tab:example-inter-case}
  \centering
  \small
  \begin{tabularx}{\linewidth}{@{}X*{2}{>{\centering\arraybackslash}p{0.17\linewidth}}@{}}
    \toprule
    \textbf{Feature} & $\boldsymbol{W_{\mathrm c}}$ & $\boldsymbol{W_{\mathrm d}}$ \\
    \midrule
    Case arrivals & Two & Two \\
    Case completions & Two & Two \\
    Active cases & Two & Two \\
    Events per active case & Two & Two \\
    Mean within-case event gap & $g$ & $\operatorname{mean}(g,G)$ \\
    \bottomrule
  \end{tabularx}
\end{table*}

\begin{figure*}[!t]
  \centering
  \begin{tikzpicture}[
  kairo/base, x=1cm, y=1cm,
  receive/.style={circle, minimum size=4.3pt, inner sep=0pt, draw=KairoTeal,
    line width=0.75pt, fill=white},
  resolve/.style={rectangle, minimum size=4.1pt, inner sep=0pt, fill=KairoTeal},
  gap/.style={draw=KairoTeal, line width=0.9pt, line cap=round, text=KairoTeal}
]
  \path[use as bounding box] (0,0) rectangle (13.08,-2.20);
  % Use the same local time positions as the resource diagram.
  \foreach \windowid/\leftx/\windowtitle/\focalend/\focalgap/\windowmean in {
    c/2.10/Earlier/2.84/g/g,
    d/7.81/Later/4.49/G/{\operatorname{mean}(g,G)}
  } {
    \begin{scope}[shift={(\leftx,0)}]
      \node[font=\footnotesize\bfseries] at (2.635,-0.17)
        {\windowtitle\ window $W_{\mathrm{\windowid}}$};
      \path[fill=KairoPaper, rounded corners=2pt] (0,-1.66) rectangle (5.27,-0.50);
      \draw[gap] (1.46,-0.91) -- node[midway, above=1pt, inner sep=1pt] {$g$} (3.56,-0.91);
      \draw[gap] (0.74,-1.43) -- node[midway, above=1pt, inner sep=1pt] {$\focalgap$}
        (\focalend,-1.43);
      \node[receive] at (1.46,-0.91) {};
      \node[resolve] at (3.56,-0.91) {};
      \node[receive] at (0.74,-1.43) {};
      \node[resolve] at (\focalend,-1.43) {};
      \path[fill=KairoTealTint, rounded corners=2pt] (0,-2.15) rectangle (5.27,-1.77);
      \node[text=KairoTeal] at (2.635,-1.96) {$\windowmean$};
    \end{scope}
  }
  \node[anchor=east] at (1.96,-0.91) {Companion};
  \node[anchor=east] at (1.96,-1.43) {Focal case};
  \node[anchor=east, font=\footnotesize\bfseries] at (1.96,-1.96) {Window mean};
\end{tikzpicture}

  \caption{The focal case's longer gap raises the window mean in $W_{\mathrm d}$. Open circles mark arrivals and filled squares mark completions.}
  \label{fig:example-inter-case}
\end{figure*}


\subsection{Scaling and Reducing the Vectors}\label{sec:approach-reduction}

Feature columns can measure counts, proportions, or durations. A column with a large numerical range can dominate clustering distances. \Cref{def:cluster-representation} specifies how Kairo rescales columns and optionally projects observations to fewer dimensions. Case identifiers and timestamps remain attached to each row as metadata.

\begin{Definition}[Scaling and projection]\label{def:cluster-representation}
Fix a perspective $p$. The representation used for clustering and its projection diagnostics are defined as follows.
\begin{itemize}[leftmargin=*,itemsep=.2em,topsep=.3em]
  \item \textbf{Column scaling.} Write $x_{pi}(o)$ for coordinate $i$ of $x_p(o)$. A scaled coordinate is
  \[
     \widetilde{x}_{pi}(o)=\frac{x_{pi}(o)-\mu_{pi}}{\sigma_{pi}},
  \]
  where $\mu_{pi}$ and $\sigma_{pi}$ are the column mean and sample standard deviation for standard score scaling, or the column minimum and range for min-max scaling. A constant column selected for scaling maps to zero. Excluded columns keep their original values.
  \item \textbf{Principal component projection.} Let $\overline{x}_p$ be the mean of the scaled vectors and let $U_\ell$ contain the leading $\ell$ orthonormal covariance eigenvectors in decreasing order of eigenvalue. With principal component analysis (PCA), $z_p(o)=U_\ell^\top(\widetilde{x}_p(o)-\overline{x}_p)$. Without projection, $z_p(o)=\widetilde{x}_p(o)$.
  \item \textbf{Projection diagnostics.} Let $\xi_1,\ldots,\xi_{d_p}$ be the covariance eigenvalues in the same order. For positive total variance, the retained fraction is $\sum_{i=1}^{\ell}\xi_i/\sum_{i=1}^{d_p}\xi_i$. The mean squared reconstruction error is
  \[
     \frac{1}{|O_p|}\sum_{o\in O_p}
       \bigl\|\widetilde{x}_p(o)-\overline{x}_p-U_\ell z_p(o)\bigr\|^2.
  \]
\end{itemize}
\end{Definition}

The activity indicators can retain their original scale. PCA can compress related coordinates, such as activity presence and frequency~\cite{hotelling1933components}. Retained variance and reconstruction error help the analyst choose the number of components. A discarded direction can still carry evidence of change.

\subsection{Learning States}\label{sec:approach-states}

Clustering groups similar observations into recurring patterns. \Cref{def:state-model} specifies how labels describe the groups and how their order gives a trajectory.

\begin{Definition}[State model and trajectories]\label{def:state-model}
Fix a perspective $p$ and the clustering representation of \Cref{def:cluster-representation}.
\begin{itemize}[leftmargin=*,itemsep=.2em,topsep=.3em]
  \item \textbf{Labels and states.} A fitted model assigns a label $s_p(o)$ to each $o\in O_p$ using its representation $z_p(o)$. The observed label set is $S_p=\{s_p(o):o\in O_p\}$. The state for $s\in S_p$ consists of $\{o\in O_p:s_p(o)=s\}$.
  \item \textbf{Trajectory.} Applying the labels along a case trace $\tau_{c_0}$ gives an intra-case trajectory. Applying them along $\sigma_p$ gives the trajectory of a window-based perspective.
  \item \textbf{State visit.} A visit is a maximal contiguous subsequence with a common label. It begins at the first observation's anchor and ends at the next visit's start. The final visit ends at the last event's timestamp for a case, or at the end of the final feature window for a window-based trajectory.
\end{itemize}
\end{Definition}

Kairo fits the model once per perspective for the analysis presented here. The resulting labels are comparable across periods, and each fit uses the complete log. \Cref{def:clustering-rules} gives the assignment rules.

\begin{Definition}[Clustering rules]\label{def:clustering-rules}
Let $z=z_p(o)$ be the representation of an observation $o\in O_p$.
\begin{itemize}[leftmargin=*,itemsep=.2em,topsep=.3em]
  \item \textbf{Distances.} For vectors $z$ and $v$, the Euclidean distance is $d_{\mathrm E}(z,v)=\|z-v\|$. For nonzero vectors, the cosine distance is $d_{\mathrm C}(z,v)=1-\langle z,v\rangle/(\|z\|\|v\|)$. The function $d_{\mathrm{cl}}$ denotes the distance selected for a method that permits a choice.
  \item \textbf{$k$-means.} The model assigns each observation to its closest center in Euclidean distance. Fitting alternates assignment with moving each occupied center to the mean of its assigned vectors~\cite{lloyd1982quantization}. For $k$ centers $\mu_1,\ldots,\mu_k$, the objective is to minimize
  \[
      \sum_{o\in O_p}\|z_p(o)-\mu_{s_p(o)}\|^2 .
  \]
  \item \textbf{Self-organizing map.} A self-organizing map (SOM) puts a prototype $\mu_g$ at each grid cell $g\in G$. Its label for $z$ is a winning cell $g^\ast$ that minimizes $d_{\mathrm{cl}}(z,\mu_g)$. In training, a prototype takes the step $\beta H(g,g^\ast)(z-\mu_g)$, where $\beta$ is the learning rate and $H$ weights cells by their proximity to the winner~\cite{kohonen1982maps}.
  \item \textbf{Density-based clustering.} Density-based spatial clustering of applications with noise (DBSCAN) uses a radius $\varepsilon$ and a minimum count $m$. A vector is a core if the sum of multiplicities of distinct vectors within distance $\varepsilon$, including itself, is at least $m$. Connected components of neighboring core vectors form clusters. A non-core vector in a core neighborhood joins a cluster of such a core. The remaining vectors receive a noise label~\cite{ester1996density}. Kairo maps each label back to the original observations with that vector.
\end{itemize}
\end{Definition}

The number of centers controls how finely $k$-means can split the observations. The SOM grid supports inspection of neighboring prototypes. For DBSCAN, the dashboard shows the ordered distances at which each observation's neighborhood reaches the chosen count, including the observation itself. A fit that finds a cluster can include noise in its state shares.

\autoref{fig:approach-clustering} uses a schematic resource-feature cloud to illustrate the assignment rules. The cloud has variations around the Frontline and Specialist resolution patterns from the running example, along with an isolated observation.

\begin{figure}[htbp]
  \centering
  \subfloat[$k$-means\label{fig:approach-clustering-kmeans}]{%
    \FigurePlaceholder[0.28]{TikZ sketch of the shared points, fitted centroids, and nearest-centroid assignments.}}
  \hfill
  \subfloat[Self-organizing map\label{fig:approach-clustering-som}]{%
    \FigurePlaceholder[0.28]{TikZ sketch of the shared points linked to their winning cells on a grid with neighboring prototypes.}}
  \hfill
  \subfloat[DBSCAN\label{fig:approach-clustering-dbscan}]{%
    \FigurePlaceholder[0.28]{TikZ sketch of core neighborhoods and border points. Mark the isolated observation as noise.}}
  \caption{State assignment on a shared schematic feature cloud. Each point represents a resource feature window. The panels show the nearest centroid for $k$-means, the winning cell for a self-organizing map, and density-connected groups and noise for DBSCAN.}
  \label{fig:approach-clustering}
\end{figure}


An analyst can interpret a state by examining the original feature values of its members. Plots of between-state distances use cluster centers for $k$-means, cell prototypes for a SOM, and means of distinct core vectors for DBSCAN. These representatives support visualization. The drift signal uses observation counts by state.

\paragraph{Inspecting the fitted states.}
\Cref{def:cluster-diagnostics} gives the measures used to check how the model organizes observations.

\begin{Definition}[State diagnostics]\label{def:cluster-diagnostics}
\leavevmode
\begin{itemize}[leftmargin=*,itemsep=.2em,topsep=.3em]
  \item \textbf{SOM quantization error.} Average $d_{\mathrm{cl}}(z_p(o),\mu_{s_p(o)})$ over $o\in O_p$.
  \item \textbf{SOM occupancy entropy.} Let $f_g$ denote the share of observations assigned to cell $g$ of grid $G$. The normalized entropy is $-\sum_{g:f_g>0}f_g\log f_g/\log|G|$ if $|G|>1$, and zero for a single-cell grid.
  \item \textbf{Silhouette.} For observation $o$, let $a_o$ be its mean distance to other members of its state and $b_o$ its least mean distance to the members of another state. The score averages $(b_o-a_o)/\max(a_o,b_o)$ over the evaluated observations, with a contribution of zero for a singleton state or a zero denominator~\cite{rousseeuw1987silhouettes}.
  \item \textbf{Calinski-Harabasz score.} Let $n$ be the number of evaluated observations, $k$ the number of their occupied states, and $n_s$ the number in state $s$. Let $\bar z_s$ be the mean vector of state $s$ and $\bar z$ the overall mean. When $2\leq k<n$ and within-state variation is positive, the score is
  \[
     \frac{\sum_s n_s\|\bar z_s-\bar z\|^2/(k-1)}
          {\sum_o\|z_p(o)-\bar z_{s_p(o)}\|^2/(n-k)}.
  \]
  The numerator measures separation of the state means and the denominator measures variation within states~\cite{calinski1974dendrite}.
\end{itemize}
The silhouette and Calinski-Harabasz calculations exclude DBSCAN noise.
\end{Definition}

These measures describe the fitted representation. \Cref{sec:assessment} evaluates how well the resulting drift signals locate annotated changes.

\subsection{Measuring Changes in State Frequencies}\label{sec:approach-signal}

The detector builds a time series from \emph{signal windows}. Each signal window collects the labels of observations whose time anchors fall inside it. \Cref{def:state-distributions} gives the grouping rule and the associated proportions. Short signal windows improve time resolution but may contain too few observations to estimate proportions reliably. \autoref{fig:approach-signal-windows} illustrates how resource feature windows contribute labels to the signal.

\begin{Definition}[Signal windows and state distributions]\label{def:state-distributions}
For perspective $p$, use the observation sequence $\sigma_p$ and time anchors $t_p$ from \Cref{def:feature-representation}.
\begin{itemize}[leftmargin=*,itemsep=.2em,topsep=.3em]
  \item \textbf{Signal-window membership.} Apply \Cref{def:calendar-grid} to $\sigma_p$ with anchor map $t_p$ and width $\Delta>0$. For a window-based perspective require $\Delta\geq\delta_p$. Write $V_1,\ldots,V_m$ for the intervals with a nonempty observation subsequence, in order of start time.
  \item \textbf{Observed state shares.} Let $O_{p,j}=\operatorname{set}((\sigma_p)_{V_j})$ and $N_j=|O_{p,j}|$. Using the labels from \Cref{def:state-model}, the distribution on $S_p$ in interval $V_j$ is
  \[
     P_j(s)=\frac{|\{o\in O_{p,j}:s_p(o)=s\}|}{N_j},
     \qquad s\in S_p.
  \]
\end{itemize}
\end{Definition}

A feature window with an empty event subsequence can still contribute a state label because it has a feature row. A signal interval with no observation anchors has no state distribution.

The reference determines what counts as recent or typical behavior. \Cref{def:drift-reference} describes a local comparison with the previous populated window, a smoothed comparison with several recent windows, and a fixed retrospective comparison with the whole log.

\begin{Definition}[Reference distribution]\label{def:drift-reference}
For a populated signal window $V_j$ from \Cref{def:state-distributions}, let $I_j$ be the reference index set.
\begin{itemize}[leftmargin=*,itemsep=.2em,topsep=.3em]
  \item \textbf{Previous window.} Use $I_j=\{j-1\}$ if $j>1$.
  \item \textbf{Recent windows.} For an integer lookback $h\geq1$, use $I_j=\{j-h,\ldots,j-1\}$ if $j>h$.
  \item \textbf{Whole log.} Use all populated indices, $I_j=\{1,\ldots,m\}$, including the current index.
\end{itemize}
When the index set is available, the reference is the unweighted mean $Q_j=|I_j|^{-1}\sum_{i\in I_j}P_i$.
\end{Definition}

\Cref{def:drift-score} compares an observed distribution with its reference. Each comparison places the resulting score at the start of the window it describes.

\begin{Definition}[Divergence and drift score]\label{def:drift-score}
For distributions $P,Q$ on $S_p$, let $M=(P+Q)/2$ and let $\epsilon>0$ be a fixed smoothing constant.
\begin{itemize}[leftmargin=*,itemsep=.2em,topsep=.3em]
  \item \textbf{Smoothed Kullback-Leibler term.} Define
  \[
    K_\epsilon(P,Q)=
      \sum_{s\in S_p}(P(s)+\epsilon)\log\frac{P(s)+\epsilon}{Q(s)+\epsilon}.
  \]
  \item \textbf{Distribution comparisons.} With square roots taken componentwise, the available functions are
  \[
    \begin{aligned}
      D_{\mathrm{KL}}(P,Q)&=K_\epsilon(P,Q),\\
      D_{\mathrm{JS}}(P,Q)&=\bigl(K_\epsilon(P,M)+K_\epsilon(Q,M)\bigr)/2,\\
      D_{\mathrm{TV}}(P,Q)&=\tfrac12\sum_{s\in S_p}|P(s)-Q(s)|,\\
      D_{\mathrm H}(P,Q)&=\|\sqrt P-\sqrt Q\|/\sqrt{2}.
    \end{aligned}
  \]
  \item \textbf{Drift score.} For a chosen function $D$ and an available reference $Q_j$, the score at $b(V_j)$ is $\lambda_j=D(P_j,Q_j)$. The score is undefined when the reference is unavailable.
\end{itemize}
\end{Definition}

These comparisons are smoothed Kullback-Leibler divergence, Jensen-Shannon divergence, total variation, and Hellinger distance. The smoothing keeps the logarithmic terms finite when a state is absent in a window. Jensen-Shannon averages differences to a common midpoint~\cite{lin1991divergence}. Total variation sums absolute changes in state shares. Hellinger uses the Euclidean distance between the vectors of square-root shares, with normalization.

\begin{Example}[From state labels to a drift score]
\label{ex:drift-score}
Consider the resource perspective of the service-desk log in \Cref{ex:service-desk}. Suppose the fitted model assigns $W_{\mathrm a},W_{\mathrm b},W_{\mathrm c}$ to state $\mathsf R$, in which Frontline resolves both cases. It assigns $W_{\mathrm d},W_{\mathrm e},W_{\mathrm f}$ to state $\mathsf S$, in which Specialist resolves the focal case. Choose a signal width twice the feature width, with $b(W_{\mathrm a})$ on a signal-window boundary. The grouping rule in \Cref{def:state-distributions} yields
\[
  V_{\mathrm A}=W_{\mathrm a}\cup W_{\mathrm b},\qquad
  V_{\mathrm B}=W_{\mathrm c}\cup W_{\mathrm d},\qquad
  V_{\mathrm C}=W_{\mathrm e}\cup W_{\mathrm f},
\]
as depicted in \autoref{fig:approach-signal-windows}.

The state distributions $P_{\mathrm A}$, $P_{\mathrm B}$, and $P_{\mathrm C}$ count feature-window labels within the respective signal windows and divide by the number of labels. All labels in $V_{\mathrm A}$ are $\mathsf R$, the labels in $V_{\mathrm B}$ are evenly split, and all labels in $V_{\mathrm C}$ are $\mathsf S$ (\autoref{fig:approach-state-distributions}). The middle window mixes the earlier and later patterns because it spans the change.

To score $V_{\mathrm C}$, the previous-window reference is $Q_{\mathrm C}^{\mathrm{prev}}=P_{\mathrm B}$. The recent-window reference that uses both earlier signal windows is $Q_{\mathrm C}^{\mathrm{recent}}=\operatorname{mean}(P_{\mathrm A},P_{\mathrm B})$. Its share of $\mathsf S$ is smaller than in $P_{\mathrm B}$ because $V_{\mathrm A}$ contains only $\mathsf R$. The whole-log reference also includes $P_{\mathrm C}$. Its mean is $P_{\mathrm B}$ here because the exclusive $\mathsf R$ and $\mathsf S$ distributions balance each other (\autoref{fig:approach-reference-distributions}).

For these states, total variation reduces to the absolute difference in the share of $\mathsf S$. Write $\alpha=P_{\mathrm B}(\mathsf S)$ for the middle window's share of this state. This gives equal previous-window scores at $b(V_{\mathrm B})$ and $b(V_{\mathrm C})$.
\[
  D_{\mathrm{TV}}(P_{\mathrm B},P_{\mathrm A})
  =D_{\mathrm{TV}}(P_{\mathrm C},P_{\mathrm B})
  =\alpha.
\]
The recent-window reference gives $V_{\mathrm C}$ a higher score because its reference contains less $\mathsf S$. The whole-log mean gives $V_{\mathrm B}$ a zero score and the same positive score to both outer windows. \autoref{fig:approach-drift-scores} places each value at the start of the window it describes. Local scores are undefined until the required earlier signal windows are available.
\end{Example}

\begin{figure*}[!t]
  \centering
  \begin{tikzpicture}[
  kairo/base,
  x=0.93cm,
  row label/.style={anchor=east, align=right, inner sep=0pt, text=KairoMuted},
  window edge/.style={draw=KairoLine, line width=0.45pt},
  signal edge/.style={draw=KairoMuted, line width=0.55pt}
]
  \node[row label] at (-0.21,1.66) {Signal window};
  \node[row label] at (-0.21,1.04) {Feature window};
  \node[row label] at (-0.21,0.45) {State at start};

  \foreach \fromX/\toX/\windowName/\shade in
    {0/4/A/KairoPaper,4/8/B/KairoTealTint,8/12/C/KairoPaper} {
    \fill[fill=\shade] (\fromX,1.41) rectangle (\toX,1.91);
    \draw[signal edge] (\fromX,1.41) rectangle (\toX,1.91);
    \node at ({(\fromX+\toX)/2},1.66) {$V_{\mathrm{\windowName}}$};
  }

  \foreach \fromX/\windowName/\state/\shade/\ink in
    {0/a/R/KairoBlueTint/KairoBlue,2/b/R/KairoBlueTint/KairoBlue,
     4/c/R/KairoBlueTint/KairoBlue,6/d/S/KairoAmberTint/KairoAmber,
     8/e/S/KairoAmberTint/KairoAmber,10/f/S/KairoAmberTint/KairoAmber} {
    \fill[fill=\shade] (\fromX,0.80) rectangle ++(2,0.48);
    \draw[window edge] (\fromX,0.80) rectangle ++(2,0.48);
    \node at ({\fromX+1},1.04) {$W_{\mathrm{\windowName}}$};
    \draw[window edge] (\fromX,0.80) -- (\fromX,0.45);
    \fill[fill=\ink] (\fromX,0.45) circle[radius=1.8pt];
    \node[anchor=west, inner sep=2.5pt, text=\ink] at (\fromX,0.45)
      {$\mathsf{\state}$};
  }
  \foreach \boundary in {0,4,8,12} {
    \draw[signal edge] (\boundary,1.41) -- (\boundary,0.80);
  }
  \draw[KairoMuted, -{Latex[length=1.65mm]}, line width=0.45pt]
    (0,0.08) -- (12,0.08) node[above left, inner sep=1.5pt] {Time};
\end{tikzpicture}

  \caption{Signal windows in the service-desk example. In the middle signal window, Frontline resolves both cases in the earlier feature window and Specialist resolves the focal case in the later one.}
  \label{fig:approach-signal-windows}
\end{figure*}

\begin{figure*}[!t]
  \centering
  \begin{tikzpicture}[kairo/base, x=1cm, y=1cm]
  \foreach \fromX/\windowName/\barDescription in {
    0/A/{Frontline only},
    4.50/B/{Equal shares},
    9.00/C/{Specialist resolves the focal case}%
  } {
    \node[anchor=west, inner sep=0pt] at (\fromX,1.32) {$P_{\mathrm{\windowName}}$};
    \node[anchor=east, inner sep=0pt, text=KairoMuted] at ({\fromX+4.10},1.32)
      {Window $V_{\mathrm{\windowName}}$};
    \node[anchor=north, inner sep=0pt, text=KairoMuted] at ({\fromX+2.05},0.53)
      {\barDescription};
  }

  \foreach \fromX/\toX/\state/\shade/\ink in {
    0/4.10/R/KairoBlueTint/KairoBlue,
    4.50/6.55/R/KairoBlueTint/KairoBlue,
    6.55/8.60/S/KairoAmberTint/KairoAmber,
    9.00/13.10/S/KairoAmberTint/KairoAmber%
  } {
    \fill[fill=\shade] (\fromX,0.65) rectangle (\toX,1.10);
    \draw[draw=\ink, line width=0.5pt] (\fromX,0.65) rectangle (\toX,1.10);
    \node[text=\ink] at ({(\fromX+\toX)/2},0.875) {$\mathsf{\state}$};
  }
\end{tikzpicture}

  \caption{Observed state distributions in the signal windows. The first and last windows contain only $\mathsf R$ and $\mathsf S$, respectively. The middle window contains equal shares.}
  \label{fig:approach-state-distributions}
\end{figure*}

\begin{figure*}[!t]
  \centering
  \begin{tikzpicture}[kairo/base, x=1cm, y=1cm]
  \foreach \fromX/\barTitle/\barSource/\ink in {
    0/{Current window}/{$P_{\mathrm C}$}/KairoInk,
    3.52/{Previous window}/{$P_{\mathrm B}$}/KairoMuted,
    6.88/{Recent mean}/{$\operatorname{mean}(P_{\mathrm A},P_{\mathrm B})$}/KairoTeal,
    10.24/{Whole-log mean}/{Equals $P_{\mathrm B}$ here}/KairoMuted%
  } {
    \node[inner sep=0pt, text=\ink, font=\footnotesize\bfseries]
      at ({\fromX+1.435},1.32) {\barTitle};
    \node[anchor=north, inner sep=0pt, text=KairoMuted]
      at ({\fromX+1.435},0.53) {\barSource};
  }
  \draw[KairoLine, line width=0.5pt] (3.19,0.22) -- (3.19,1.45);

  \foreach \fromX/\toX/\state/\shade/\ink in {
    0/2.87/S/KairoAmberTint/KairoAmber,
    3.52/4.955/R/KairoBlueTint/KairoBlue,
    4.955/6.39/S/KairoAmberTint/KairoAmber,
    6.88/9.0325/R/KairoBlueTint/KairoBlue,
    9.0325/9.75/S/KairoAmberTint/KairoAmber,
    10.24/11.675/R/KairoBlueTint/KairoBlue,
    11.675/13.11/S/KairoAmberTint/KairoAmber%
  } {
    \fill[fill=\shade] (\fromX,0.65) rectangle (\toX,1.10);
    \draw[draw=\ink, line width=0.5pt] (\fromX,0.65) rectangle (\toX,1.10);
    \node[text=\ink] at ({(\fromX+\toX)/2},0.875) {$\mathsf{\state}$};
  }
\end{tikzpicture}

  \caption{Reference distributions for scoring $V_{\mathrm C}$. The recent-window reference has less $\mathsf S$ than the previous-window and whole-log references in this example.}
  \label{fig:approach-reference-distributions}
\end{figure*}

\begin{figure*}[!t]
  \centering
  \begin{tikzpicture}[
  kairo/base,
  x=1cm,
  y=1cm,
  grid line/.style={draw=KairoLine, line width=0.45pt},
  value line/.style={line width=0.91pt, line cap=round}
]
  \node[anchor=east, inner sep=0pt, text=KairoMuted] at (0.08,2.32)
    {Score};
  \foreach \level/\tickLabel in {0.38/$0$,1.10/$\alpha$,1.46/$\tfrac32\alpha$} {
    \node[anchor=east, inner sep=0pt, text=KairoMuted] at (0.08,\level) {\tickLabel};
  }

  \foreach \fromX/\panelTitle/\ink in {
    0.24/{Previous window}/KairoBlue,
    4.65/{Recent mean}/KairoTeal,
    9.06/{Whole-log mean}/KairoMuted%
  } {
    \node[font=\footnotesize\bfseries, inner sep=0pt, text=\ink]
      at ({\fromX+2.02},2.32) {\panelTitle};
    \fill[KairoPaper] (\fromX,0.38) rectangle ++(4.04,1.63);
    \foreach \level in {1.10,1.46}
      \draw[grid line] (\fromX,\level) -- ++(4.04,0);
    \draw[KairoMuted, line width=0.5pt] (\fromX,0.38) -- ++(4.04,0);
    \foreach \offset/\windowName in {0.60/A,2.02/B,3.44/C} {
      \draw[KairoMuted, line width=0.45pt]
        ({\fromX+\offset},0.38) -- ++(0,-0.06);
      \node[anchor=north, text=KairoMuted, inner sep=0pt]
        at ({\fromX+\offset},0.25) {$b(V_{\mathrm{\windowName}})$};
    }
  }

  \draw[value line, KairoBlue] (2.26,1.10) -- (3.68,1.10);
  \foreach \x in {2.26,3.68}
    \fill[KairoBlue] (\x,1.10) circle[radius=2.0pt];
  \fill[KairoTeal] (8.09,1.46) circle[radius=2.0pt];
  \draw[value line, KairoMuted] (9.66,1.10) -- (11.08,0.38) -- (12.50,1.10);
  \foreach \x/\y in {9.66/1.10,11.08/0.38,12.50/1.10}
    \fill[KairoMuted] (\x,\y) circle[radius=2.0pt];
\end{tikzpicture}

  \caption{Illustrative total-variation scores. The scale uses the middle window's share $\alpha$ of state $\mathsf S$. Missing points indicate that the reference needs more history. The whole-log reference gives the middle window a score of zero.}
  \label{fig:approach-drift-scores}
\end{figure*}


A high score indicates an interval to inspect. It can reflect a change in process behavior or ordinary variation, particularly when a window has low support or covers only part of the monitored period. Different perspectives can respond to the same change. \Cref{sec:assessment-config} defines the benchmark's alarm rule for turning these scores into estimated change times.

\section{Implementation}\label{sec:implementation}

\subsection{Analysis Library}

Kairo is a Python library that implements the analysis in \autoref{sec:approach}. Feature extraction, scaling, dimensionality reduction, clustering, and drift scoring are available as separate stages. This lets an analyst inspect intermediate results, choose features that the log supports, and reuse the same state and drift analysis with different features. The library links observations back to their events or calendar windows, so the states and scores can be examined in the context of the original log. The benchmark uses the same analysis library as the graphical interface.

\subsection{Graphical Interface}

The browser-based dashboard guides analysts from log upload to drift inspection. It accepts eXtensible Event Stream (XES) and comma-separated values (CSV) files and lets the analyst map the case, activity, and timestamp fields. It also supports optional resource and case attributes. When start and completion timestamps are available, the analyst can derive duration for the resource features.

The analyst opens a perspective and chooses its feature groups, scaling, and clustering settings. The dashboard displays feature tables and plots of explained variance and reconstruction error to help choose whether and how far to reduce the vectors. After clustering, the interface shows state frequencies and trajectories. The analyst can follow individual cases in the intra-case view or follow the log through successive feature windows in the resource and inter-case views. Comparing state frequencies across selected periods helps reveal which patterns became more or less common.

The drift view places the divergence signal alongside state distributions and observation counts for the same periods. The analyst can adjust the threshold and inspect the support behind a peak, as illustrated in \autoref{fig:implementation-dashboard}. An optional Copilot lets the analyst ask questions about selected summaries and plots. Its interpretations are intended for review against the displayed evidence. The detector and the benchmark run independently of the Copilot.

\begin{figure}[H]
  \centering
  \FigurePlaceholder{Screenshot of the Kairo drift view showing the divergence signal, adjustable threshold, state distributions, and observation counts.}
  \caption{Dashboard view for investigating a candidate drift. The analyst can compare a high score with the state shares and observation count behind it.}
  \label{fig:implementation-dashboard}
\end{figure}

\section{Assessment}
\label{sec:assessment}

We assess how well Kairo localizes annotated changes in complete event logs. The event logs, detector configurations, and scoring protocol provide a common basis for comparison with implementations in \softwarepm{}. \Cref{sec:assessment-results} addresses \ref{rq:comparison} with results pooled across logs. \Cref{sec:assessment-scenarios} answers \ref{rq:scenarios} through selected change scenarios. \Cref{sec:assessment-sensitivity} addresses \ref{rq:margin} by repeating the scoring with different matching margins.

\subsection{Event Logs}
\label{sec:assessment-logs}

We selected 23 event logs with 156,383 cases, 1,086,668 events, and 37 annotated drift episodes. \autoref{tab:assessment-logs} lists all inputs. The selection covers several kinds of change and uses each selected log in full. The Padella collection injects changes into logs from real processes~\cite{padella2026dataset}. Our sample uses eight files from five base processes, including two controls with no injected drift. The nine synthetic Rheon logs from the thesis include sudden and gradual changes to activity order, assignments, workload, waiting time, and case composition~\cite{kairo-thesis}. The six historical logs from the archive of Adams et~al. contain synthetic control-flow changes~\cite{adams2023experimental}. They include a Ceravolo scenario with 15\% trace noise~\cite{tavares2019synthetic}.

\begin{table}[H]
\caption{Event logs used in the comparison. Simultaneous annotations with identical time bounds count as one episode.}
\label{tab:assessment-logs}
\centering
\small
\begin{adjustwidth}{-\extralength}{0cm}
\begin{tabularx}{\fulllength}{@{}lrrr>{\raggedright\arraybackslash}X@{}}
\toprule
\textbf{Log} & \textbf{Cases} & \textbf{Events} & \textbf{Episodes} & \textbf{Annotated changes} \\
\midrule
Bose & 6,000 & 58,838 & 4 & Sudden control-flow changes \\
Ceravolo RIO & 1,000 & 11,395 & 1 & Sudden control-flow change \\
Ceravolo RIO (15\% noise) & 1,000 & 10,554 & 1 & Sudden control-flow change \\
Ceravolo RP & 1,000 & 11,009 & 1 & Sudden control-flow change \\
Ostovar Loop & 2,999 & 100,471 & 2 & Sudden changes involving a loop \\
Ostovar Parallel to Sequence & 2,999 & 112,011 & 2 & Sudden changes involving parallel and sequential execution \\
\midrule
Padella 04 & 8,049 & 51,059 & 1 & Gradual speedup in loan applications \\
Padella 05 & 10,500 & 52,308 & 1 & Sudden control-flow change in domestic declarations \\
Padella 06 & 32,429 & 212,721 & 1 & Sudden speedup in bank account closures \\
Padella 07 & 6,886 & 36,796 & 0 & Request for payment control \\
Padella 08 & 6,886 & 36,796 & 1 & Sudden control-flow change in requests for payment \\
Padella 10 & 10,500 & 56,437 & 1 & Sudden slowdown in domestic declarations \\
Padella 12 & 8,049 & 51,059 & 0 & Loan application control \\
Padella 15 & 8,000 & 39,493 & 1 & Gradual control-flow change in purchase orders \\
\midrule
Rheon Single 03 & 3,254 & 12,992 & 1 & Sudden control-flow change \\
Rheon Single 11 & 8,896 & 24,543 & 1 & Sudden resource reassignment \\
Rheon Single 15 & 9,042 & 72,684 & 1 & Gradual change in waiting time \\
Rheon Multi 02 & 3,161 & 20,514 & 4 & Sudden control-flow changes \\
Rheon Multi 03 & 9,523 & 23,852 & 3 & Sudden changes in the regional case mix \\
Rheon Multi 07 & 3,132 & 15,660 & 3 & Changes in workload, resource pool size, and assignments \\
Rheon Mixed 06 & 4,796 & 23,980 & 1 & Coincident workload and waiting-time changes \\
Rheon Mixed 07 & 2,860 & 14,300 & 2 & Changes in workload and case arrival rates \\
Rheon Mixed 13 & 5,422 & 37,196 & 4 & Changes in control flow, workload, and case arrival rates \\
\midrule
\textbf{Total} & \textbf{156,383} & \textbf{1,086,668} & \textbf{37} & \\
\bottomrule
\end{tabularx}
\end{adjustwidth}
\end{table}

We expressed all annotations in calendar time. For Rheon we used generator timestamps. For Padella we used the start times of the cases identified in the published annotations. For Bose and Ostovar we used the start of the first post-change trace in the original archive order. The Ceravolo archive stores numeric case identifiers in lexicographic order, so we used the start of case 500, the first case in the changed half by numeric generator identifier. We preserved physical trace boundaries when the Bose archive reused a case identifier.

\subsection{Detectors and Configurations}
\label{sec:assessment-config}

The comparison covers nine Kairo combinations of perspective and clustering method and 18 \softwarepm{} implementations (\autoref{tab:assessment-methods}). We call a specific configuration of an implementation a \emph{profile}. Core profile settings used observable input properties, such as log length, event and case volume, and available fields. The annotated change times were used for scoring after detection. The core runs used Kairo 0.1.0 at revision \texttt{6b1bd462} and \softwarepm{} at revision \texttt{79b6ce1b}, with revision \texttt{765435bc} for the corrected Bose detector. The benchmark archive records the complete settings for every run. Coverage depended on data and computational requirements. The \softwarepm{} resource-assignment profile ran only when the log recorded resources. Kairo recorded 30 unsupported runs due to missing resources and eight because the number of distinct intra-case feature rows exceeded the DBSCAN limit of 10,000. In total, 616 \softwarepm{} runs and 376 Kairo runs completed.

\begin{table}[H]
\caption{Implementations and profiles included in the core comparison. DFG denotes a directly-follows graph. PELT denotes pruned exact linear time, ADWIN denotes adaptive windowing, and MMD denotes maximum mean discrepancy.}
\label{tab:assessment-methods}
\centering
\small
\begin{adjustwidth}{-\extralength}{0cm}
\begin{tabularx}{\fulllength}{@{}l>{\raggedright\arraybackslash}p{0.19\fulllength}>{\raggedright\arraybackslash}X@{}}
\toprule
\textbf{Library} & \textbf{Implementation} & \textbf{Signal and tested profiles} \\
\midrule
Kairo & State distributions & Intra-case, resource, and inter-case states with K-means, a self-organizing map (SOM), or DBSCAN. Each combination uses previous-window and recent-window profiles. \\
\midrule
\softwarepm{} & Bose relation count & Permutation test on counts of outgoing eventually-follows relations. \\
& DFG frequency & Adjacent-window comparison of relation frequencies using permutations or a wider likelihood-ratio test. \\
& Event relations & Fixed or adaptive windows of event-level control-flow relations. \\
& Local completeness & New or disappearing directly-follows relations relative to a locally complete window. \\
& Relation invariance & Boundaries of stable periods for activity relations. \\
& Graph metrics & PELT on measures of successive DFGs. \\
& Run distribution & Fixed or adaptive windows of partially ordered execution runs. \\
& Trace variant frequency & Permutation tests on exact activity sequences with regular and wide windows. \\
& Earth Mover distance & Window comparisons with activity sequences alone or with timing features. \\
& Calendar indicators & PELT on basic or enriched indicator sets, and ADWIN on basic indicators. \\
& CV-based indicators & PELT on enriched indicators with smoothing selected using the coefficient of variation. \\
& Conformance & Alignment fitness monitored with fixed-model Page-Hinkley or ADWIN with relearning. \\
& Conformance profile & PELT on windowed alignment fitness and precision. \\
& Declare profiles & PELT on temporal constraint profiles. \\
& Version clustering & Clustering and segmentation of window similarity profiles based on DFGs. \\
& Cluster frequency & Comparison of trace cluster frequencies after training K-means on an initial set of traces. \\
& Multivariate features & Permutation MMD on flow and timing features or on activity-resource assignments. \\
& Supervised similarity & A patch classifier on window similarity images trained on independent generated logs. \\
\bottomrule
\end{tabularx}
\end{adjustwidth}
\end{table}

\paragraph{Kairo profiles.}
The state models were fitted to each complete log. The benchmark applied min-max scaling to intra-case vectors and standard score scaling to resource and inter-case vectors, followed by principal component analysis with up to five components when applicable. It used observed fields and feature size limits to choose the vectors. The base signal window targeted one fiftieth of the log's completion-time span, rounded to an available duration. Resource and inter-case feature windows were chosen from lengths that divide the base signal window, taking event or case density into account. The previous-window profile used Jensen-Shannon divergence and the 90th percentile of eligible scores as its threshold. The recent-window profile used total variation relative to the mean of the preceding three populated windows, twice the base signal window width, and the 95th percentile. Both profiles required full coverage and at least 25\% of the median number of state observations. \Cref{def:assessment-alarms} specifies how these settings produce retrospective alarms from the scores.

\begin{Definition}[Retrospective alarm rule]\label{def:assessment-alarms}
Let $t_{\min}$ and $t_{\max}$ be the first and last completion timestamps in the input log, and let $\widetilde N$ be the median of the populated-window counts $N_j$ from \cref{def:state-distributions}. For a chosen support fraction $\rho$, the set of eligible indices is
\[
  J=\{j:\lambda_j\text{ is finite},\ V_j\subseteq[t_{\min},t_{\max}),\
               N_j\geq\rho\widetilde N\}.
\]
If $J$ is nonempty, compute the empirical quantile at the chosen level from the sequence $(\lambda_j)_{j\in J}$, including repeated scores. The threshold $\theta$ is that quantile with a floor of $10^{-12}$. A window qualifies if $j\in J$ and $\lambda_j\geq\theta$. A maximal run of qualifying windows whose calendar boundaries touch produces an alarm at the start of its first window. An empty $J$ produces no alarms.
\end{Definition}

\paragraph{\softwarepm{} profiles.}
We used the recorded \softwarepm{} source revisions and retained each implementation's own alarms. Trace window sizes generally scaled with case count, and calendar indicators used bins sized by log duration. When usable service timestamps were missing, timing profiles used gaps between completions. Version clustering used 64 trace windows with a minimum cluster size of four and a segmentation penalty of three. The supervised similarity profile used a local classifier fitted to 24 independent synthetic control-flow logs. For the conformance profile, precision was calculated with the memory-efficient exact Dijkstra solver. On Ostovar Parallel to Sequence, exact alignment took too long with the model first mined from the initial traces. The conformance profiles used a sequential reference based on the medoid of the first 30 completed traces. Detection began after that training prefix, and the adaptive profile could relearn from one completed case following an alarm.

Before scoring the Bose detector, we corrected two implementation issues. The detector had used case identifiers to order traces and counted incoming eventually-follows relations. We ordered cases by their first event and used outgoing counts, as specified by the Relation Type Count definition~\cite{bose2011handling}. A check with permuted case identifiers confirmed that chronology was preserved. On the historical Bose log, the relation-count correction preserved three of four matches at the primary tolerance and reduced unmatched alarms from 17 to 12.

\subsection{Scoring Protocol}
\label{sec:assessment-scoring}

We used a common calendar-time margin equal to 2.5\% of the log's observation span, from the earliest recorded start or completion to the last completion. We repeated scoring with margins of 1\% and 5\%. Coincident annotations with identical bounds form one episode. \Cref{def:assessment-matching} applies the same localization rule to every detector, regardless of the type of change.

\begin{Definition}[Temporal matching]\label{def:assessment-matching}
An annotated episode $i$ has bounds $[l_i,r_i]$, with $l_i=r_i$ for a point change. The distance from an alarm timestamp $t$ to this episode is
\[
  d(t,[l_i,r_i])=\max\{l_i-t,\ 0,\ t-r_i\}.
\]
For a given margin $\gamma$, an alarm and an episode can be paired when their distance is at most $\gamma$. Among one-to-one matchings of feasible pairs, we maximize the number of matches and then minimize the sum of their distances.
\end{Definition}

Each match contributes a true positive ($TP$), each unmatched alarm a false positive ($FP$), and each unmatched episode a false negative ($FN$). Unmatched alarms include repeat alarms and changes that may have occurred outside the annotated injections. We sum $TP$, $FP$, and $FN$ across logs before computing micro precision $TP/(TP+FP)$, micro recall $TP/(TP+FN)$, and micro F1 $2TP/(2TP+FP+FN)$. The no-injection controls therefore contribute unmatched alarms to the pooled counts. We score the estimated change times. Measuring reporting delay would require the confirmation times, which the benchmark records where available.

\subsection{Comparative Performance Across Logs}
\label{sec:assessment-results}

For \ref{rq:comparison}, \autoref{tab:assessment-main-results} shows the three highest micro F1 profiles from each library among the profiles that completed on all 23 logs. At the primary margin, Kairo's previous-window intra-case K-means profile matched 21 of 37 episodes, while \softwarepm{} version clustering matched 15. Kairo raised 72 unmatched alarms and version clustering raised 10. Nine of Kairo's unmatched alarms and two of version clustering's came from the control logs. Version clustering therefore led in precision and micro F1, and the Kairo profile led in recall. Both profiles matched five episodes in the historical logs. On Padella, Kairo matched four with 27 unmatched alarms and version clustering matched three with seven unmatched alarms. On Rheon, Kairo matched 12 with 30 unmatched alarms and version clustering matched seven without an unmatched alarm.

\begin{table}[htbp]
\caption{Highest micro F1 profiles with full log coverage at the primary tolerance of 2.5\%. Rows are ranked within each library.}
\label{tab:assessment-main-results}
\centering
\small
\begin{tabularx}{\textwidth}{@{}l>{\raggedright\arraybackslash}Xrrrrrr@{}}
\toprule
\textbf{Library} & \textbf{Profile} & \textbf{TP} & \textbf{FP} & \textbf{FN} & \textbf{Precision} & \textbf{Recall} & \textbf{F1} \\
\midrule
Kairo & Intra-case K-means, previous & 21 & 72 & 16 & 0.226 & 0.568 & \textbf{0.323} \\
& Inter-case SOM, previous & 22 & 80 & 15 & 0.216 & 0.595 & 0.317 \\
& Intra-case SOM, previous & 20 & 73 & 17 & 0.215 & 0.541 & 0.308 \\
\softwarepm{} & Version clustering & 15 & 10 & 22 & 0.600 & 0.405 & \textbf{0.484} \\
& Declare profiles, PELT & 10 & 9 & 27 & 0.526 & 0.270 & 0.357 \\
& Trace variant frequency, wide & 20 & 66 & 17 & 0.233 & 0.541 & 0.325 \\
\bottomrule
\end{tabularx}
\end{table}

The coverage restriction matters for interpreting the pooled comparison. Kairo's recent-window resource K-means profile reached micro F1 of 0.370 on its 18 supported logs. On the same 18 logs, version clustering reached 0.440. We report such profiles separately because the omitted logs alter the set of episodes and opportunities for unmatched alarms.

\subsection{Detection by Change Scenario}
\label{sec:assessment-scenarios}

For \ref{rq:scenarios}, we compare profiles in selected scenarios at the primary matching margin. The examples show which annotated changes the perspective-specific state signals can capture under the tested settings.

\paragraph{Control flow.}
On the three Padella control-flow logs, Kairo's recent-window intra-case SOM and \softwarepm{} version clustering each matched all three injections and each raised two unmatched alarms. On Rheon Multi 02, Kairo's previous-window intra-case SOM matched all four sudden control-flow changes and raised one unmatched alarm. The wide trace variant frequency profile matched all four with no unmatched alarms. Both DFG frequency profiles matched all three Ceravolo injections with no unmatched alarms, including the log with trace noise.

\paragraph{Assignments and timing.}
All three recent-window Kairo resource profiles matched the sudden reassignment in Rheon Single 11 with no unmatched alarms. The \softwarepm{} resource MMD and enriched indicator PELT profiles also matched it and each raised one unmatched alarm. Kairo's recent-window inter-case SOM matched the gradual waiting-time change in Rheon Single 15 with no unmatched alarms. The trained trace cluster frequency profile matched it and raised one unmatched alarm. \autoref{fig:assessment-rheon-signals} pairs the Kairo resource and inter-case SOM signals for these two changes. On the three Padella timing logs, previous-window resource SOM and fixed-model Page-Hinkley conformance each matched two injections and raised nine unmatched alarms. Both missed the speedup in bank account closures.

\begin{figure}[H]
  \centering
  % Data for the side-by-side assessment panels.
% The source paths below are relative to assessments/concept-drift/results/.
% Calendar coordinates are normalized from the first window start to the last window end.
% Scores keep their original scale, with coordinates rounded to six decimal places.

% rheon_single_11/dataset.json and rheon_single_11/kairo/resource-som/recent.json.
\def\RheonReassignmentPath{%
  \draw[rheon/series] plot[mark=*,mark size=1.05pt] coordinates {
    (0.093750,0.240741) (0.125000,0.092593) (0.156250,0.074074) (0.187500,0.148148)
    (0.218750,0.074074) (0.250000,0.092593) (0.281250,0.240741) (0.312500,0.111111)
    (0.343750,0.259259) (0.375000,0.259259) (0.406250,0.777778) (0.437500,0.759259)
    (0.468750,0.444444) (0.500000,0.314815) (0.531250,0.222222) (0.562500,0.111111)
    (0.593750,0.111111) (0.625000,0.277778) (0.656250,0.111111) (0.687500,0.240741)
    (0.718750,0.074074) (0.750000,0.129630) (0.781250,0.240741) (0.812500,0.055556)
    (0.843750,0.222222) (0.875000,0.148148) (0.906250,0.203704) (0.937500,0.092593)
    (0.968750,0.444444)
  };%
}
\def\RheonReassignmentEdge{(0.968750,0.444444)}
\def\RheonReassignmentAlarm{(0.406250,0.777778)}
\def\RheonReassignmentThreshold{0.649074}
\def\RheonReassignmentOnset{0.414687}
\def\RheonReassignmentTicks{0.026042/Apr 2020, 0.343750/Jun, 0.661458/Aug, 0.979167/Oct}

% rheon_single_15/dataset.json and rheon_single_15/kairo/inter-som/recent.json.
\def\RheonWaitingPath{%
  \draw[rheon/series] plot[mark=*,mark size=1.05pt] coordinates {
    (0.100000,0.392857) (0.133333,0.273810) (0.166667,0.464286) (0.200000,0.190476)
    (0.233333,0.047619) (0.266667,0.107143) (0.300000,0.095238) (0.333333,0.071429)
    (0.366667,0.166667) (0.400000,0.607143) (0.433333,0.869048) (0.466667,0.785714)
    (0.500000,0.511905) (0.533333,0.369048) (0.566667,0.285714) (0.600000,0.142857)
    (0.633333,0.047619) (0.666667,0.154762) (0.700000,0.642857) (0.733333,0.369048)
    (0.766667,0.250000) (0.800000,0.535714) (0.833333,0.190476) (0.866667,0.345238)
    (0.900000,0.750000) (0.933333,0.761905) (0.966667,0.531746)
  };%
}
\def\RheonWaitingEdge{(0.966667,0.531746)}
\def\RheonWaitingAlarm{(0.433333,0.869048)}
\def\RheonWaitingThreshold{0.779762}
\def\RheonWaitingOnset{0.364286}
\def\RheonWaitingEnd{0.468571}
\def\RheonWaitingTicks{0.016667/2015, 0.233929/2016, 0.451786/2017, 0.669048/2018, 0.886310/2019}

  \begin{tikzpicture}[kairo/base,x=1cm,y=1cm,
      rheon/grid/.style={draw=KairoLine,line width=.35pt},
      rheon/axis/.style={draw=KairoMuted,line width=.5pt},
      rheon/series/.style={draw=KairoBlue,line width=.8pt,line join=round},
      rheon/threshold/.style={draw=KairoMuted,line width=.6pt,dash pattern=on 3pt off 2pt},
      rheon/truth/.style={draw=KairoAmber,line width=.65pt},
      rheon/alarm/.style={rectangle,rotate=45,inner sep=0pt,minimum size=5pt,
                          fill=KairoTeal,draw=white,line width=.45pt},
      rheon/tick/.style={inner sep=0pt,text=KairoMuted}]
    \node[anchor=west,font=\small\bfseries,inner sep=0pt] at (.83,2.61)
      {(a) Resource reassignment};
    \node[anchor=west,font=\small\bfseries,inner sep=0pt] at (7.39,2.61)
      {(b) Gradual waiting-time change};
    \node[rotate=90,inner sep=0pt,text=KairoMuted] at (.02,1.14) {Total variation};

    \begin{scope}[shift={(.83,0)},x=5.89cm,y=2.28cm]
      \fill[KairoPaper] (0,0) rectangle (1,1);
      \draw[rheon/grid] (0,.5) -- (1,.5) (0,1) -- (1,1);
      \draw[rheon/truth] (\RheonReassignmentOnset,0) -- (\RheonReassignmentOnset,1);
      \draw[rheon/threshold] (0,\RheonReassignmentThreshold) -- (1,\RheonReassignmentThreshold);
      \RheonReassignmentPath
      \filldraw[draw=KairoBlue,fill=white,line width=.55pt] \RheonReassignmentEdge circle[radius=1.5pt];
      \node[rheon/alarm] at \RheonReassignmentAlarm {};
      \draw[rheon/axis] (0,1) -- (0,0) -- (1,0);
      \foreach \y/\value in {0/0,.5/0.5,1/1} {
        \draw[rheon/axis] (0,\y) -- ++(-2.3pt,0);
        \node[rheon/tick,anchor=east,xshift=-5pt] at (0,\y) {\value};
      }
      \foreach \x/\date in \RheonReassignmentTicks {
        \draw[rheon/axis] (\x,0) -- ++(0,-2.3pt);
        \node[rheon/tick,anchor=north,yshift=-5pt] at (\x,0) {\date};
      }
    \end{scope}

    \begin{scope}[shift={(7.39,0)},x=5.89cm,y=2.28cm]
      \fill[KairoPaper] (0,0) rectangle (1,1);
      \fill[KairoAmberTint] (\RheonWaitingOnset,0) rectangle (\RheonWaitingEnd,1);
      \draw[rheon/grid] (0,.5) -- (1,.5) (0,1) -- (1,1);
      \draw[rheon/truth] (\RheonWaitingOnset,0) -- (\RheonWaitingOnset,1)
        (\RheonWaitingEnd,0) -- (\RheonWaitingEnd,1);
      \draw[rheon/threshold] (0,\RheonWaitingThreshold) -- (1,\RheonWaitingThreshold);
      \RheonWaitingPath
      \filldraw[draw=KairoBlue,fill=white,line width=.55pt] \RheonWaitingEdge circle[radius=1.5pt];
      \node[rheon/alarm] at \RheonWaitingAlarm {};
      \draw[rheon/axis] (0,1) -- (0,0) -- (1,0);
      \foreach \x/\date in \RheonWaitingTicks {
        \draw[rheon/axis] (\x,0) -- ++(0,-2.3pt);
        \node[rheon/tick,anchor=north,yshift=-5pt] at (\x,0) {\date};
      }
    \end{scope}

    \draw[rheon/series] (1.02,-.82) -- (1.49,-.82);
    \fill[KairoBlue] (1.255,-.82) circle[radius=1.05pt];
    \node[anchor=west,inner sep=0pt] at (1.61,-.82) {Window score};
    \draw[rheon/threshold] (4.06,-.82) -- (4.53,-.82);
    \node[anchor=west,inner sep=0pt] at (4.65,-.82) {Alarm threshold};
    \fill[KairoAmberTint] (7.36,-.93) rectangle (7.83,-.71);
    \draw[rheon/truth] (7.595,-.93) -- (7.595,-.71);
    \node[anchor=west,inner sep=0pt] at (7.95,-.82) {Annotated drift};
    \node[rheon/alarm] at (10.85,-.82) {};
    \node[anchor=west,inner sep=0pt] at (11.04,-.82) {Estimated alarm};
  \end{tikzpicture}
  \caption{Recent-window Kairo SOM scores for Rheon Single 11 (a) and Rheon Single 15 (b). The annotated change is a point for resource reassignment and an interval for gradual waiting-time change. The diamond marks the alarm at the start of its signal window. Hollow circles mark finite scores from windows excluded from the alarm rule.}
  \label{fig:assessment-rheon-signals}
\end{figure}


\paragraph{Case composition.}
On Rheon Multi 03, the dominant region changed three times and returned to region 4 at the last change. Kairo's inter-case SOM with a previous-window reference matched all three shifts and raised two unmatched alarms. The recent-window version matched two with no unmatched alarms. The core \softwarepm{} PELT indicators and flow-and-timing MMD profiles matched none. We then ran a diagnostic with the case attributes supplied explicitly. PELT used region shares and average case amounts, and MMD used region as a categorical feature and amount as a numeric feature. Each detected all three shifts with zero unmatched alarms, including at the 1\% margin. This diagnostic evaluates an attribute choice made after looking at the scenario. It is reported separately from the core ranking.

\subsection{Sensitivity to the Matching Margin}
\label{sec:assessment-sensitivity}

For \ref{rq:margin}, we took the profiles that led each library at the primary margin in \autoref{tab:assessment-main-results} and scored them at both alternative margins. At 1\%, micro F1 was 0.231 for Kairo and 0.419 for version clustering. At 5\%, their scores were 0.400 and 0.516. Version clustering retained the higher micro F1 at both margins.

With the alarms fixed, the margin changes which alarms can match the annotations. For example, the wide DFG likelihood-ratio profile placed two boundaries in Ostovar Loop about 26 hours after their annotated case starts. Both boundaries matched only at the widest margin.

\section{Discussion}
\label{sec:discussion}

At the main matching margin, Kairo's strongest overall profile matched more annotated changes than the strongest \softwarepm{} profile. The \softwarepm{} profile achieved a better balance between precision and recall at every margin tested. Additional alarms give an analyst more intervals to investigate and increase the effort required to monitor a process. Both profiles also raised alarms on logs without injected changes. Those alarms may reflect ordinary fluctuations or other changes outside the benchmark annotations, and their operational importance would need separate investigation.

The examples help explain when a state representation is useful. A resource profile surfaced a reassignment, and an inter-case profile surfaced a gradual increase in waiting time. Activity sequences alone provide little direct evidence about either change. Where the injected change did affect activity sequences, methods that compare sequences or activity relations performed well. In particular, Kairo and version clustering gave the same aggregate results on the selected Padella control-flow logs. The case-composition experiments illustrate another part of the picture. Once the comparison methods were given the relevant case attributes explicitly, they located the region changes cleanly. The choice of measurements can matter as much as the choice of the detector.

Separate state descriptions give analysts different views of the same episode. An analyst can inspect case progress, resource activity, and overall load around the same time and compare the measurements behind the states that became more or less common. Co-occurring signals can guide that investigation. For instance, a reassignment may change workload and waiting time together. The assessment compared the timing of alarms with annotations. Whether the interface helps analysts explain a change, or trace a causal link between perspectives, was outside that evaluation.

The benchmark was retrospective. Kairo learned the states and the relative alarm threshold from the full log, and the strongest profiles were identified on the logs used for scoring. Its performance therefore describes retrospective localization on the selected benchmark. The alarm threshold marks large changes relative to the rest of the same log. A stable process can still produce peaks, and a log with several large changes can provide a demanding reference for smaller ones. The log families were chosen for coverage of different behaviors and change types. The comparison cannot establish how often the same results would occur in operational logs or with prospective monitoring.

\section{Conclusion and Future Work}
\label{sec:conclusion}

Building on Deren's thesis~\cite{kairo-thesis}, this paper formalizes how Kairo learns states from case progress, resource activity, and conditions across cases, then uses changes in their frequency to locate drift. The library and dashboard let analysts inspect the states and the measurements behind their signals. The archived benchmark documents how these signals located selected resource and timing changes. It also quantifies the trade-off between finding more episodes and raising more unmatched alarms.

Future work should test fixed profiles on independent logs using the information available at each alarm time. These tests could examine whether thresholds calibrated on stable periods and rules that suppress repeat alerts reduce review effort. A study with analysts could assess whether the dashboard helps them identify what changed.


\authorcontributions{B.D. developed and implemented the state-based method and created the Rheon generator. A.B. advised the thesis and conducted the comparative benchmark.}
\funding{Funding details are pending author confirmation.}
\dataavailability{The \href{https://github.com/fit-alessandro-berti/kairo-thesis/tree/a2a6d60d757e0bb04804152162f853c4a42ca087}{project repository} provides the Kairo library and Rheon generator. Its \href{https://github.com/fit-alessandro-berti/kairo-thesis/tree/a2a6d60d757e0bb04804152162f853c4a42ca087/assessments/concept-drift}{benchmark archive} provides reproduction instructions, run configurations, drift annotations, and results. The run records identify the source of each event log.}
\acknowledgments{OpenAI Codex assisted in drafting and organizing the manuscript and checking its descriptions against the existing thesis, software, and benchmark outputs.}
\conflictsofinterest{Conflict of interest declarations are pending author confirmation.}

\reftitle{References}
\bibliography{bib/core,bib/related-work,bib/method,bib/assessment}

\end{document}
